% VI49-DETAILED-CHALLENGE-SOLUTIONS
% VI49-DETAILED-CHALLENGE-01
\subsection*{Challenge VI.49.1}
\begin{solution}
If one member \(U_0\) of a cover is the whole space, define \(h:C^p\to C^{p-1}\) by inserting the distinguished index \(0\) into each ordered index tuple, with the sign needed to restore order.  Expanding the alternating differential shows
\[
\delta h+h\delta=\mathrm{id}
\]
in positive degrees; in degree zero the same identity gives exactness of the augmentation.  This is the contracting homotopy used after localizing the principal-cover complex.
\end{solution}

% VI49-DETAILED-CHALLENGE-02
\subsection*{Challenge VI.49.2}
\begin{solution}
A standard nonseparated example is obtained by gluing two copies of \(\mathbb A^2\) along \(\mathbb A^2\setminus\{0\}\) by the identity.  The two copies are affine opens, but their intersection is the punctured affine plane, which is not affine.  Thus finite affine intersections need not remain affine without separatedness, and an affine cover may fail to be Leray.
\end{solution}

% VI49-DETAILED-CHALLENGE-03
\subsection*{Challenge VI.49.3}
\begin{solution}
The \(n+1\) standard affine charts of \(\mathbb P^n\) form a Leray cover for \(\mathcal O(d)\), so cohomology vanishes above degree \(n\).  Multivariable Laurent-monomial analysis predicts that line bundles have no intermediate cohomology and that only \(H^0\) for sufficiently nonnegative \(d\) and the top group for sufficiently negative \(d\) survive.  The \(\mathbb P^1\) missing-exponent interval becomes a lattice-region calculation.
\end{solution}

% VI49-DETAILED-CHALLENGE-04
\subsection*{Challenge VI.49.4}
\begin{solution}
Starting with
\[
0\to\mathcal F\to\mathcal I^0\to\mathcal Q^1\to0
\]
and using acyclicity of \(\mathcal I^0\),
\[
H^q(\mathcal F)\cong H^{q-1}(\mathcal Q^1)\quad(q\ge2).
\]
Repeat through the flabby resolution until
\[
H^q(\mathcal F)\cong H^1(\mathcal Q^{q-1}).
\]
Thus all higher groups can be reduced to a first obstruction group.
\end{solution}

% VI49-DETAILED-CHALLENGE-05
\subsection*{Challenge VI.49.5}
\begin{solution}
Basic affine vanishing is local-to-global through affine/Leray covers and requires no positivity; it applies to quasi-coherent sheaves on affine schemes and supplies cover-length bounds on separated schemes.  Serre vanishing is asymptotic and projective: for coherent \(\mathcal F\) and ample twisting sheaf, \(H^q(X,\mathcal F(n))=0\) for \(q>0\) once \(n\gg0\).  Its proof uses finite generation and projective graded-module machinery beyond this volume.
\end{solution}

